% !TeX root = ../papers/nc-p3.tex
\section{二次曲面射影の二つの母線が与える中間辺と正則延長の障害}
\label{sec:quadric-biform-window}
$Q=\mathbb P^1_s\times\mathbb P^1_t$ とし、二項二重次数切断空間を $S_{ab}$、
射影を与える四空間を $V\subset S_{22}$ と書く。
各母線上で $V$ の制限が二項二次式の全空間を張ると仮定する。
\begin{theorem}\label{thm:quadric-ruling-edge-obstruction}
二つの商空間 $A=S_{23}/V S_{01}$, $B=S_{32}/V S_{10}$ は四次元である。
$y\in\mathbb P(S_{11})$ の両側に $m\in\mathbb P(A^*)$, $n\in\mathbb P(B^*)$ を置き、
\[
            m(yh)=0\ (h\in S_{12}),\qquad n(yk)=0\ (k\in S_{21})
\]
を二辺の式とする。これらを関係空間とする四辺有限な標準Koszul左右AS正則延長は存在しない。
二十次射影の正規化上、長さ三繊維の各母線への一次切断の制限が二独立なら、
その繊維の張る平面 $y$ に対し、上の商でこの二つの中間座標を取れる
（当該繊維への $S_{23},S_{32}$ の制限が全射の場合）。
\end{theorem}
\begin{proof}
$t$ を底とする評価列とその捻りを取る：
\[
 0\longrightarrow\mathcal O(-6)\longrightarrow V\otimes\mathcal O
       \longrightarrow S_{20}\otimes\mathcal O(2)\longrightarrow0 .
\]
全射性は母線仮定で、核の次数は行列式で決まる。
$k=-1,0,1$ の捻りの接続写像は $S_{2,k+2}/V S_{0k}\simeq H^1(\mathcal O(k-6))$ を与える
（負次数は零）。双対の $t$ の一次式による写像は二項式の通常の積である。
$s$ と $t$ を交換しても成立する。
特に $A^*$ は $t$ の三次式、$B^*$ は $s$ の三次式であり、
$(S_{22}/V)^*$ の二つの同定から可逆写像 $T:S_t^4\to S_s^4$ を得る。ここで $S_u^j$ は $u$ の二項 $j$ 次式空間。
$y=\sum y_{ij}s_it_j$ と置く。$h\mapsto m(t_j s_i h)$ をまず $S_{22}/V$ で読み、
次いで $S_{12}^*=S_s^5$ に移すと、両辺は
\[
          \sum_{i,j} y_{ij}s_i T(t_jm)=0,\qquad
          \sum_{i,j} y_{ij}t_j T^{-1}(s_i n)=0
\]
となる。この転置された係数は命題\ref{prop:trisecant-common-quadric}の標準中間二辺と同じである。
二つの有理正規四次曲線が異なるなら定理\ref{thm:trisecant-middle-simple-obstruction}が
第四文字の選び方によらず正則延長を排除する。一致するなら $T$ は二項変数の基底変更の
四次対称冪のスカラー倍であり、変数を合わせると第一式は $m\sum y_{ij}s_i s_j=0$ となる。
交代な $y$ には全 $m$ が対応するので有限辺にならない。

射影繊維 $Y$ 上では $V$ の像は単元で生成される一直線である。$S_{23}$ の繊維上の値を
$V S_{01}$ の値で割ると、長さ三から二を引いた一次元となり、その双対が $m$ を与える。
$y$ は繊維で消えるので全ての $yh$ が消される。逆向きも同じである。
\end{proof}

\begin{proposition}\label{prop:quadric-rational-frobenius-counterexample}
前節の有理四空間では上の二辺を含む四辺有限・標準14641 Frobenius窓が実在し、
反対切断階数は十六である。その四周期二次代数には頂点次元 $(3,3,2,2)$ の単純表現があり、
全次数の正則性に失敗する。
\end{proposition}
\begin{proof}
残りの二辺を明記する。$x,m$ 辺は
\[
       R_{xm}=\ker\{ V\otimes A\longrightarrow S_{45}/(\operatorname{Sym}^2 V) S_{01}\}
\]
とし、$n,x$ 辺は $s,t$ を交換した核を転置する。商のため代表の選択によらない。
長さ三繊維から来る上の $x,m$ はこの式も満たす。実際 $V$ を繊維で $x$ のスカラー列にし、
積の等式に $m$ を定めた一次汎関数を施せばよい。
\path{computations/verify_quadric_biform_window.py} は四空間からこれらを計算する。
各辺の六関係の両射影で一般平面の代わりに指定平面を用い、十五最大小行列式が全ての平面四次式を張ることを
素数1009で証明するので、標数零でも有限である。四始点の隣接交叉次元は
次数三、四、五で $4,1,0$。頂テンソルは巡回で定数倍となり、各切断の階数 $4,6,4$ は
双対の全積対合の完全性を与える。これらの階数・交叉と反対切断は有理数上で計算する。

単純表現も具体的に構成できる。$D_t=S_{24}/V S_{02}$, $D_s=S_{42}/V S_{20}$ は三次元で、
頂点を順に $D_s,D_t^*,k^2,k^2$ とする。$y$ の四矢は二行二列の行列単位。
$m$ 位置の矢 $a\in A$ は、$p\in D_t^*$ を $(p(t_1a),-p(t_0a))$ に送る。
$n$ 位置は同じ $s$ に関する写像の転置である。二つの中央の関係は交代積で消える。
$x$ の四行列の三十六成分に残り二辺の線形条件を課すと、厳密計算で解空間は四次元、
その一元は四行列の共通核が零で像の総和が全三空間である。
各母線での全射性は三次小行列式の二項六次式に五次式を掛けて十一次全空間を張ることで検証する。
前の評価列を $k=2$ にも用いると $D_t^*=S_t^2$ であり、$m$ の作用は非零 $p$ について
$t_0p,t_1p$ の独立性から次の二空間を張る。転置側でも各非零二ベクトルに対し像が三空間を張る。
中間の行列単位と $x$ の共通核・像条件を併せると非零部分表現は全ての頂点を満たす。
よって単純であり、定理\ref{thm:generic-corner-small-simples}の必要条件に反する。
この排除は指定した連続二辺に対するもので、同じ反対対への他の中間座標はなお調べる必要がある。
\end{proof}
